%%%%%%%%%%%%%%%%%%%%%%%%%%%%%%%%%%%%%%%%%%%%%%%%%%%%%%%%%%%%%%%%%%
\section{Entries of the approximate Jacobian}\label{app:jacobian}
%%%%%%%%%%%%%%%%%%%%%%%%%%%%%%%%%%%%%%%%%%%%%%%%%%%%%%%%%%%%%%%%%%

This appendix gives the entries of the blocks of the approximate Jacobian $\tilde{\bb{J}}$ of Equations~\eqref{eqn:block_preconditioner} and~\eqref{eqn:fluid_block_preconditioner}, summarised in Section~\ref{sec:approximate_block_preconditioner}.
$\mathcal{F}^{\mathrm{int}}_P$ denotes the internal faces of cell $P$, $\bb{d}_f$ the vector joining the centres of the two cells sharing face $f$, and $\hat{\bb{n}}_f$ the unit face normal.

For a fluid cell $P$, the diagonal block of the approximate momentum operator
is formed from the compact diffusion contribution, the BDF2 time derivative and
a first-order upwind linearisation of the convection term,
\begin{align}
    \left[\tilde{\bb{J}}_{mm}\right]_{PP}
    =
    &-
    \sum_{f \in \mathcal{F}^{\mathrm{int}}_P}
    \nu
    \frac{|\bb{\Gamma}_{F,f}|}
         {\hat{\bb{n}}_{F,f} \cdot \bb{d}_{F,f}}
    \bb{I}
    -
    \frac{3}{2\Delta t}
    \Omega_{F,P}
    \bb{I}
    \notag \\
    &-
    \sum_{f \in \mathcal{F}^{\mathrm{int}}_P}
    \max\left(\phi^{\mathrm{rel}}_{F,f},0\right)\bb{I},
    \label{eqn:Jmm_diag}
\end{align}
while the off-diagonal block associated with a neighbouring fluid cell $Q$
across face $f$ is
\begin{equation}
    \left[\tilde{\bb{J}}_{mm}\right]_{PQ}
    =
    \nu
    \frac{|\bb{\Gamma}_{F,f}|}
         {\hat{\bb{n}}_{F,f} \cdot \bb{d}_{F,f}}
    \bb{I}
    -
    \min\left(\phi^{\mathrm{rel}}_{F,f},0\right)\bb{I}.
    \label{eqn:Jmm_offdiag}
\end{equation}
The explicit transpose part of the viscous stress is not included in $\tilde{\bb{J}}_{mm}$; the matrix-free Jacobian action accounts for it.
Here the convective contribution in the approximate Jacobian is always based on first-order upwinding, independently of the convection scheme used in the nonlinear residual (linear or limited second-order upwind interpolation in the cases considered).

The pressure-gradient coupling block is obtained by differentiating the cell-centred pressure-gradient term in the momentum equation with respect to the cell-centred pressure.
If $\bb{g}_{P,Q}$ denotes the least-squares gradient coefficient multiplying the pressure in cell $Q$ when evaluating $(\bb{\nabla}p)_P$, then
\begin{equation}
    \left[\tilde{\bb{J}}_{mp}\right]_{P Q}
    =
    -
    \Omega_{F,P}
    \bb{g}_{P,Q},
    \label{eqn:Jmp}
\end{equation}
which is a $d\times 1$ block in $d$ spatial dimensions.
The continuity block is formed from the derivative of the discrete pressure
equation with respect to the cell-centred velocity. For an internal face $f$
shared by cells $P$ and $Q$, using the interpolation weight $w_f$ for the value
from cell $P$,
\begin{align}
    \left[\tilde{\bb{J}}_{pm}\right]_{PP}
    &=
    -
    \sum_{f \in \mathcal{F}^{\mathrm{int}}_P}
    w_f \bb{\Gamma}_{F,f}^{T},
    \label{eqn:Jpm_diag}
    \\
    \left[\tilde{\bb{J}}_{pm}\right]_{PQ}
    &=
    -
    \left(1-w_f\right)\bb{\Gamma}_{F,f}^{T},
    \label{eqn:Jpm_offdiag}
\end{align}
which are $1\times d$ blocks.

The pressure block $\tilde{J}_{pp}$ is obtained from the pressure-stabilisation
term. For a compact Laplacian-type representation of the stabilisation,
\begin{align}
    \left[\tilde{J}_{pp}\right]_{PP}
    &=
    -
    \sum_{f \in \mathcal{F}^{\mathrm{int}}_P}
    \frac{\alpha_p}{a_{P,f}}
    \frac{|\bb{\Gamma}_{F,f}|}
         {\hat{\bb{n}}_{F,f} \cdot \bb{d}_{F,f}},
    \label{eqn:Jpp_diag}
    \\
    \left[\tilde{J}_{pp}\right]_{PQ}
    &=
    \frac{\alpha_p}{a_{P,f}}
    \frac{|\bb{\Gamma}_{F,f}|}
         {\hat{\bb{n}}_{F,f} \cdot \bb{d}_{F,f}},
    \label{eqn:Jpp_offdiag}
\end{align}
where $\alpha_p/a_{P,f}$ is the face pressure-stabilisation coefficient arising
from the Rhie--Chow stabilisation in Equation~\eqref{eqn:rhie_chow}. Here
$a_{P,f}$ denotes the face-interpolated compact momentum diagonal coefficient.
In addition to this compact part, the contribution of the interpolated-gradient term in Equation~\eqref{eqn:rhie_chow}, whose stencil extends to the neighbours of the neighbours, is also inserted into $\tilde{J}_{pp}$, so that the pressure block contains the complete linear stabilisation operator.

The block $\tilde{J}_{pp}$ is therefore a compact pressure--pressure coupling
introduced by the finite-volume pressure stabilisation, rather than by a
physical pressure diffusion term.
Algebraically, it approximates the pressure coupling that would arise through the velocity block when the velocity is
eliminated from the coupled momentum--pressure system.
This connection motivates its use in the velocity--pressure split strategy discussed below.

The block
$\tilde{\bb{J}}_{fs}
=
\left[\tilde{\bb{J}}_{ms};\tilde{\bb{J}}_{ps}\right]$
represents the dependence of the fluid momentum and pressure equations on the
mapped solid interface velocity. In these blocks, the solid interface face
velocity is approximated by the velocity of the adjacent solid cell (i.e.\ the
gradient extrapolation of Section~\ref{sec:interface_conditions} is neglected),
and the transfer weights $w_{fs}$ of
Equation~\eqref{eqn:solid_to_fluid_transfer} are applied. For an interface
face $f$ and a solid face $s$ that overlaps it, the fluid
momentum-row contribution arises from the Dirichlet boundary contribution to
the compact diffusion operator and is approximated by
\begin{equation}
    \left[\tilde{\bb{J}}_{ms}\right]_{P Q}
    =
    w_{fs}\,
    \nu
    \frac{|\bb{\Gamma}_{F,f}|}
         {\hat{\bb{n}}_{F,f} \cdot \bb{d}_{F,f}}
    \bb{I},
    \label{eqn:fluid_solid_velocity_momentum_coupling}
\end{equation}
where $P$ is the fluid cell adjacent to face $f$ and $Q$ is the solid cell
adjacent to face $s$; contributions from several faces are summed. This is a
$d\times d$ block in $d$ spatial dimensions. The
pressure-equation contribution is the derivative of the interface normal
velocity flux with respect to the mapped solid velocity,
\begin{equation}
    \left[\tilde{\bb{J}}_{ps}\right]_{P Q}
    =
    -
    w_{fs}\,
    \bb{\Gamma}_{F,f}^{T},
    \label{eqn:fluid_solid_velocity_pressure_coupling}
\end{equation}
which is a $1\times d$ block. The sign follows the pressure residual definition
in Equation~\eqref{eqn:fluid_pressure_residual}.

The solid block $\tilde{\bb{J}}_{ss}$ is approximated using a compact
diffusion-like operator combined with the velocity time-derivative and damping
terms. This extends the displacement-based approximate Jacobian of
\citet{Cardiff2026} to the present velocity formulation. For a solid cell $P$,
the diagonal block is
\begin{align}
    \left[\tilde{\bb{J}}_{ss}\right]_{PP}
    =
    &-
    \sum_{f \in \mathcal{F}^{\mathrm{int}}_P}
    \frac{2\Delta t}{3}
    \bar{K}_f
    \frac{|\bb{\Gamma}_{S,0,f}|}
         {\hat{\bb{n}}_{S,0,f} \cdot \bb{d}_{S,0,f}}
    \bb{I}
    \notag \\
    &-
    \frac{3}{2\Delta t}
    \rho_{S,0} \Omega_{0,P}
    \bb{I}
    -
    c_d \rho_{S,0} \Omega_{0,P}
    \bb{I},
    \label{eqn:Jss_diag}
\end{align}
while the off-diagonal block associated with a neighbouring solid cell $Q$
across face $f$ is
\begin{equation}
    \left[\tilde{\bb{J}}_{ss}\right]_{PQ}
    =
    \frac{2\Delta t}{3}
    \bar{K}_f
    \frac{|\bb{\Gamma}_{S,0,f}|}
         {\hat{\bb{n}}_{S,0,f} \cdot \bb{d}_{S,0,f}}
    \bb{I}.
    \label{eqn:Jss_offdiag}
\end{equation}
Here $\bar{K}$ is the stiffness of Equation~\eqref{eqn:rhie_chow_s}, and the
factor $2\Delta t/3$ is the derivative of the new-time displacement in
Equation~\eqref{eqn:solid_bdf2_displacement} with respect to the solid
velocity, which arises from expressing the compact displacement
diffusion operator $\bb{\nabla}_0\cdot(\bar{K}\bb{\nabla}_0\bb{D})$, used as an
approximation of the divergence of the stress, in terms of the solid velocity
unknown.
This approximation provides a compact first-neighbour solid block for the
approximate Jacobian, but it is not a consistent linearisation of the nonlinear
total Lagrangian constitutive response.

For a solid interface face $s$ associated with solid cell $P$, overlapped by a
fluid interface face $f$ whose adjacent fluid cell is $Q$, the pressure part of
$\tilde{\bb{J}}_{sf}
=
\left[\tilde{\bb{J}}_{sm},\tilde{\bb{J}}_{sp}\right]$
gives the derivative of the solid momentum equation with respect to the
adjacent fluid pressure. With the kinematic-pressure convention, the
contribution is approximated as
\begin{equation}
    \left[\tilde{\bb{J}}_{sp}\right]_{P Q}
    =
    w_{sf}\, \rho_F |\bb{\Gamma}_{S,s}|\, \hat{\bb{n}}_{F,f},
    \label{eqn:pressure_coupling_approximation}
\end{equation}
in which the interface pressure is approximated by the pressure of the adjacent
fluid cell, $|\bb{\Gamma}_{S,s}|$ is the solid interface face area and
$\hat{\bb{n}}_{F,f}$ is the outward fluid unit normal, used, as in the residual,
because the fluid and solid normals differ on a non-conformal curved interface.
This is a $d\times 1$ block.

When viscous tractions are transferred, the derivative of the solid interface
force with respect to the adjacent fluid velocity is approximated by a compact
normal-gradient contribution,
\begin{equation}
    \left[\tilde{\bb{J}}_{sm}\right]_{P Q}
    =
    w_{sf}\,
    \rho_F \nu
    \frac{|\bb{\Gamma}_{S,s}|}
         {\hat{\bb{n}}_{F,f} \cdot \bb{d}_{F,f}}
    \left(
    \bb{I}
    +
    \hat{\bb{n}}_{F,f} \otimes \hat{\bb{n}}_{F,f}
    \right),
    \label{eqn:viscous_coupling_approximation}
\end{equation}
which is a $d\times d$ block.
The normal--normal term $\hat{\bb{n}}_{F,f} \otimes \hat{\bb{n}}_{F,f}$ is retained from the linearisation of the full viscous stress, although the residual uses the normal-gradient form of Equation~\eqref{eqn:interface_fluid_traction}; this is one of the approximations of $\tilde{\bb{J}}$.
The pressure and viscous coupling blocks are multiplied by $\rho_F$ because the fluid equations use the kinematic pressure while the solid traction is a force per unit area.
